% =====================================================================
% Introduction (style/flow pass: reads as one coherent opening)
%
% Design unchanged. Three methods (BR, ExtT, tCop) symmetric; common
% comparison object chi_{M,c}(d); seasonal contrast Delta_M^chi(d).
% No invented results (planned work uses "will"/"is designed to").
%
% Style changes in this pass:
%   - removed run-in \paragraph{} mini-headings; flowing prose;
%   - varied the phrasing for chi(d) (no repeated "canonical comparison
%     object"): "common comparison object", "the curve chi(d)", "the
%     quantity used to compare the models", "common tail-dependence
%     curve", etc.;
%   - smoother transitions, less per-sentence over-explaining;
%   - lighter than methodology/simulation.
%
% Labels preserved: sec:introduction, sec:intro-rq (referenced by
% methodology), sec:intro-approach, sec:intro-contribution,
% sec:intro-outline, eq:intro-chi, eq:intro-delta. sec:intro-outline is
% an alias on the contribution subsection.
% Citation keys: subset of the existing .bib (no new keys).
% =====================================================================
\section{Introduction}\label{sec:introduction}

Severe wind gusts are among the most damaging weather hazards in the Netherlands, and what makes them costly is rarely the peak speed at a single location. A storm that drives one station to an extreme reading while leaving the rest of the country calm is a local nuisance; a storm that pushes many stations past their extremes on the same day is a national event. Insured losses, damage to infrastructure, and disruption to transport and the electricity network all scale with how much of the country is hit at once, so the \emph{spatial extent} of an extreme event is as relevant as its local intensity. For European windstorms this extent is increasingly studied in its own right \citep{sharkey2020}, and it is the perspective this thesis takes for Dutch wind-gust extremes.

Two rather different mechanisms produce damaging gusts here, and they differ precisely in their spatial reach. In winter, the extremes are usually generated by large extratropical cyclones whose wind fields span hundreds of kilometres, so that an extreme gust at one station tends to come with extreme gusts at many others the same day. In summer, the extremes are more often produced by convective storms and downbursts, which can be locally violent but cover only a small area \citep{pacey2021,gliksman2023european}. Seen from a single station the two look alike, because in both cases that station records a severe gust. But, seen across the network they are not alike at all, and this contrast motivates the question the thesis sets out to answer: whether extreme gusts co-occur over larger distances in winter than in summer.

Turning that question into a statistical question requires deciding what to measure. The obvious candidate, correlation, will not be sufficient. Correlation is computed from all days, and almost all days are unremarkable, so it mainly describes how ordinary windy and calm weather co-moves; two pairs of stations can share nearly the same correlation while differing sharply in how often their most extreme gusts coincide. Because the question is about those rare joint extremes, the right object has to be defined on the extremes themselves. The one used throughout this thesis is the probability that one station records an extreme gust given that a neighbour does so on the same day, taken in the limit as the threshold for ``extreme'' becomes more and more severe, and read as a function of the distance between the two stations. This is the upper tail-dependence coefficient \(\chi(d)\) \citep{coles1999}: where it decays quickly with distance, joint extremes are local. And where it decays slowly, extremes tend to strike wide areas together. The decay of \(\chi(d)\) with distance is the precise version of the loose phrase ``spatial extent'' used above.

\subsection{Research question}\label{sec:intro-rq}

The thesis is built around the question:

\begin{quote}\itshape
    Does the spatial extremal dependence of Dutch daily maximum wind-gust extremes differ between winter and summer?
\end{quote}

Concretely, it asks whether tail dependence between the \(33\) station locations of the Royal Netherlands Meteorological Institute (KNMI) decays more slowly with inter-station distance in winter (December--February) than in summer (June--August). To make ``decays more slowly'' precise, we focus on a single quantity for comparison.
Here \(i\neq j\) denotes a pair of distinct KNMI stations, \(d_{ij}\) is their inter-station distance, and \(U_i,U_j\in(0,1)\) denote wind-gust observations transformed to a common uniform marginal scale, as detailed in the methodology \ref{sec:methodology}. For a fitted model \(M\), a season \(c\in\{\mathrm W,\mathrm S\}\), and uniform-scale observations \((U_i,U_j)\) at stations a distance \(d\) apart,

\begin{equation}\label{eq:intro-chi}
    \chi_{M,c}(d)=\lim_{u\to1}\Pr_{M}\bigl(U_i>u \mid U_j>u;\ d_{ij}=d\bigr),
\end{equation}
and the winter--summer difference is the gap between two such curves,
\begin{equation}\label{eq:intro-delta}
    \Delta_M^{\chi}(d)=\chi_{M,\mathrm W}(d)-\chi_{M,\mathrm S}(d),
\end{equation}
with \(\Delta_M^{\chi}(d)>0\) meaning that model \(M\) implies stronger winter tail dependence at distance \(d\). The seasonal question is answered by examining the sign and size of \(\Delta_M^{\chi}(d)\) over the distances the network resolves. The reason \(\chi(d)\) rather than any internal model parameter plays this role is that it is defined the same way for every model considered here and measures exactly the co-occurrence of extremes that the question is about; the internal ``range'' parameters, by contrast, mean different things in different models and cannot be compared directly.

\subsection{Estimation of spatial extremal dependence}\label{sec:intro-approach}

Several models can produce a curve \ref{eq:intro-chi} from assumptions that are far from equivalent. Rather than commit to one model, the thesis treats the choice of model as part of the design and estimates the seasonal curves with three parametric methods side by side.

Before any model is fitted, the station--season margins are removed. A coastal station has a windier climate than an inland one, and any station is windier in winter than in summer, so dependence measured on the raw gust scale would confound spatial co-occurrence with these marginal differences. Each station--season series is therefore standardised to a common scale. Let \(\widehat U_{itc}=\widehat F_{ic}(X_{itc})\) for the probability-integral transform of the gust \(X_{itc}\) at station \(i\), season \(c\), day \(t\), and \(\widehat Z_{itc}=-1/\log\widehat U_{itc}\) for the unit-Fr\'echet version, so that what remains to be compared is dependence and nothing else. All three models operate on these standardised data.

The three methods are chosen to span a meaningful range of assumptions while each delivering \(\chi_{M,c}(d)\) in closed form. The first two are max-stable process models, which are the limit class for spatial maxima. First, the Brown--Resnick model \citep{brown1977,kabluchko2009}, which is a stationary model in spatial extremes, whose dependence is governed by a variogram and whose curve decays to zero at long range. Second, the extremal-\(t\) model \citep{opitz2013}, an alternative within the same class (max-stable) that adds a tail parameter and imposes a different shape of decay, including a small but non-zero level of dependence at long range. The third model is a spatial Student-\(t\) copula \citep{demarta2005}. Since Student-\(t\) copulas are a tool for modelling tail dependence, and are not constrained by max-stability. All three are treated symmetrically. Each is fitted separately in winter and summer and then converted into its curve \(\widehat\chi_{M,c}(d)\) and contrast \(\widehat\Delta_M^{\chi}(d)\). One caveat in comparing these three models, is that Student-\(t\) copula implies some positive tail dependence even when the correlation is weak, so its fitted curve is read as a parametric benchmark for the seasonal contrast rather than as evidence of physical dependence between distant storms.

Estimation and inference are kept common across the three methods so that any difference in their conclusions reflects their individual assumptions. The two max-stable models cannot be fitted by full likelihood at \(33\) stations, so they are estimated by pairwise composite likelihood over all station pairs; the copula is fitted by the matching pairwise copula likelihood, so that all three see the same pairs and the same days. Because these objectives are not ordinary likelihoods and because storm days cluster in time, standard errors are taken from a season-year block bootstrap that resamples whole winters and whole summers, leaving the within-storm and same-day structure intact. Its primary output is the bootstrap distribution of the difference curve \(\Delta_M^{\chi}(d)\). Before we apply the methods to our empirical dataset, the methods will be evaluated in a simulation study to test them on data generated from known structures on the actual station geometry, measuring whether each recovers a seasonal difference that is present and avoids signalling one that is absent.

How the three methods relate then becomes the substance of the answer. If all three imply \(\widehat\Delta_M^{\chi}(d)>0\) over the distances that matter, the evidence for a wider winter footprint is robust to the assumptions that separate them; if only one does, the finding is model-sensitive and is reported as such; and if they disagree in direction, that disagreement is itself informative, because the methods differ in identifiable ways, above all in whether they impose max-stability, so a disagreement points to the assumption responsible.

\subsection{Contribution and outline}\label{sec:intro-contribution}\label{sec:intro-outline}

This thesis develops no new model, estimator, or asymptotic theory. Its aim is a careful comparative assessment: a seasonal comparison of spatial extremal dependence in Dutch wind-gust extremes, based on roughly thirty-five years of daily maxima at \(33\) KNMI stations, together with an account of whether three parametric models with different assumptions agree on the direction and distance profile of the seasonal difference. The quantity compared, \(\chi(d)\), is directly relevant to aggregate wind risk because it governs how often many locations are hit at once, and a conclusion that survives across models inside and outside the max-stable class is more credible than one tied to a single family. Methodologically, the contribution is in the design of the comparison. Three models judged on one common curve, fitted to identically standardised data with a matched estimation strategy, given uncertainty through one shared block bootstrap, and benchmarked by simulation on the observed geometry before the empirical curves are interpreted.

The remainder follows the same path. Section~\ref{sec:literature} sets out the background the design rests on, from univariate extreme-value theory and marginal standardisation through max-stable processes, the Brown--Resnick and extremal-\(t\) models, Student-\(t\) copulas and tail dependence, composite likelihood, and the block bootstrap, to the use of simulation for comparing estimators. Section~\ref{sec:methodology} specifies the procedures in full: the marginal standardisation, the curve \(\chi_{M,c}(d)\) and its seasonal contrast, the three models with their closed-form curves, the pairwise estimation, the seasonal comparison, and the block bootstrap. Section~\ref{sec:simulation} gives the design of the simulation comparison. Section~\ref{sec:data} describes the KNMI wind-gust data, and Section~\ref{sec:results} reports the simulation outcomes and the empirical seasonal estimates from the three methods, with their cross-method comparison. Section~\ref{sec:conclusion} concludes and discusses limitations and extensions.